\subsection{Масса Джинса}
\label{sec:jeans-mass}
Звёзды рождаются из холодных облаков межзвёздного газа. Сжимает облако его собственная гравитация, а удерживает от сжатия давление газа. При каких условиях гравитация побеждает и облако начинает сжиматься ответ даёт \term{критерий Джинса}, полученный Джеймсом Джинсом в 1902 году.

Рассмотрим ту же модель, что и при оценке давления в центре звезды, \lookSecRef{sec:pressure-and-temperature}: однородное шарообразное облако массы $M$, радиуса $R$, плотности $\rho$ и температуры $T$, состоящее из частиц массы $m$.

\paragraph{Гравитационная энергия облака} Будем собирать облако из тонких сферических слоёв, добавляя их снаружи один за другим. Когда собран шар радиуса $r$ и массы $M(r) = \frac{4}{3}\pi\rho r^3$, очередной слой толщиной $dr$ имеет массу $dm = 4\pi\rho r^2\,dr$, а его потенциальная энергия в поле уже собранного шара
\begin{equation*}
    dU = -\frac{G M(r)\,dm}{r} = -\frac{16}{3}\pi^2 G \rho^2 r^4\,dr.
\end{equation*}
Проинтегрируем по всем слоям и выразим плотность через полную массу, $\rho = 3M / (4\pi R^3)$:
\begin{equation}
    U = -\frac{16}{3}\pi^2 G \rho^2 \int\limits_0^R r^4\,dr
      = -\frac{16}{15}\pi^2 G \rho^2 R^5
      = -\frac{3}{5}\frac{GM^2}{R}.
    \label{eq:sphere-grav-energy}
\end{equation}

\paragraph{Критерий Джинса} Кинетическая энергия облака~--- это энергия теплового движения его $N = M / m$ частиц, \lookSecRef{sec:mkt}. Чтобы не путать её с температурой, обозначим кинетическую энергию $K$:
\begin{equation*}
    K = \frac{3}{2} N k T = \frac{3}{2}\frac{M}{m}kT.
\end{equation*}
По теореме о вириале, \lookSecRef{sec:virial}, облако находится в равновесии, если $2K = -U$. Если же тепловой энергии не хватает, $2K < -U$, гравитация побеждает и облако сжимается:
\begin{gather*}
    3\frac{M}{m}kT < \frac{3}{5}\frac{GM^2}{R},\\
    M > \frac{5kT}{Gm}R.
\end{gather*}
В это условие радиус входит вместе с массой. Исключим его с помощью плотности, $R = \left(3M / 4\pi\rho\right)^{1/3}$:
\begin{gather*}
    M > \frac{5kT}{Gm}\left(\frac{3M}{4\pi\rho}\right)^{1/3},\\
    M^{2/3} > \frac{5kT}{Gm}\left(\frac{3}{4\pi\rho}\right)^{1/3},\\
    M > \left(\frac{5kT}{Gm}\right)^{3/2}\left(\frac{3}{4\pi\rho}\right)^{1/2}.
\end{gather*}
Правая часть называется \term{массой Джинса}:
\begin{equation}
    M_J = \left(\frac{5kT}{Gm}\right)^{3/2}\left(\frac{3}{4\pi\rho}\right)^{1/2},
    \label{eq:jeans-mass}
\end{equation}
здесь $T$ и $\rho$~--- температура и плотность облака, $m$~--- масса его частиц. Облако массы больше $M_J$ не может удерживаться давлением и сжимается под действием собственной гравитации. Масса Джинса растёт с температурой, $M_J \propto T^{3/2}$, и падает с плотностью, $M_J \propto \rho^{-1/2}$: чем холоднее и плотнее газ, тем меньше массы достаточно для сжатия.

То же условие можно записать и через размер облака. Подставим в $M > 5kTR / (Gm)$ массу $M = \frac{4}{3}\pi\rho R^3$:
\begin{gather*}
    \frac{4}{3}\pi\rho R^3 > \frac{5kT}{Gm}R,\\
    R > \sqrt{\frac{15 kT}{4\pi G m \rho}}.
\end{gather*}
Правая часть называется \term{радиусом Джинса}:
\begin{equation}
    R_J = \sqrt{\frac{15 kT}{4\pi G m \rho}}.
    \label{eq:jeans-radius}
\end{equation}
Облако данной плотности и температуры сжимается, если его размер больше $R_J$, и остаётся в равновесии, если меньше.

\paragraph{Смысл радиуса Джинса} Давление противостоит гравитации не мгновенно: возмущение давления распространяется по газу со скоростью звука $c_s$, которая по порядку величины равна тепловой скорости частиц, $c_s \sim \sqrt{kT / m}$. Гравитация же сжимает облако за время свободного падения $t_\text{св}$. Оценим его: частица на краю облака испытывает ускорение $g = GM / R^2 = \frac{4}{3}\pi G \rho R$ и, если считать его постоянным, проходит путь $R$ за время
\begin{equation*}
    t_\text{св} \simeq \sqrt{\frac{2R}{g}} = \sqrt{\frac{3}{2\pi G \rho}} \sim \frac{1}{\sqrt{G\rho}}.
\end{equation*}
Заметим, что время свободного падения зависит только от плотности и не зависит от размера облака. Радиус Джинса \eqref{eq:jeans-radius} с точностью до числового множителя есть путь, который звук проходит за время свободного падения:
\begin{equation*}
    R_J \sim \sqrt{\frac{kT}{m}} \cdot \frac{1}{\sqrt{G\rho}} \sim c_s t_\text{св}.
\end{equation*}
В облаке меньше $R_J$ давление успевает выровняться быстрее, чем гравитация заметно сожмёт облако, и сжатие останавливается. В облаке больше $R_J$ давление не успевает отреагировать, и сжатие ничто не сдерживает. Строгий расчёт, учитывающий распространение волн в газе, даёт те же зависимости $M_J$ и $R_J$ от температуры и плотности, отличаясь лишь числовыми множителями порядка единицы, поэтому формулы \eqref{eq:jeans-mass} и \eqref{eq:jeans-radius} следует понимать как оценки.

\begin{wrapfigure}[14]{r}{0.48\tw}
    \centering
    \vspace{-1pc}
    \tikzsetnextfilename{jeans-mass}
    \begin{tikzpicture}
        \footnotesize
        % Масса Джинса для молекулярного водорода, m = 2 m_p:
        % M_J / M_sun = 23.1 (T / 10 К)^{3/2} (n / 10^9 м^{-3})^{-1/2}
        % Степенной закон на логарифмических осях есть прямая, поэтому
        % каждую кривую задаём двумя крайними точками, n = 1e6 и 1e12 м^-3
        \begin{loglogaxis}[
            width = 0.43\tw,
            height = 5.2cm,
            clip = false,
            xmin = 1e6, xmax = 1e12,
            ymin = 0.1, ymax = 1e5,
            xlabel = {$n,~\text{м}^{-3}$},
            ylabel = {$M_J / M_\odot$},
            xtick = {1e6, 1e8, 1e10, 1e12},
            ytick = {0.1, 10, 1e3, 1e5},
        ]
            \addplot [black] coordinates {(1e6, 730.4) (1e12, 0.7304)}
                node [pos = 1, anchor = west] {$10~\text{К}$};
            \addplot [black] coordinates {(1e6, 3795) (1e12, 3.795)}
                node [pos = 1, anchor = west] {$30~\text{К}$};
            \addplot [black] coordinates {(1e6, 23100) (1e12, 23.1)}
                node [pos = 1, anchor = west] {$100~\text{К}$};
            \addplot [dashes] coordinates {(1e6, 1) (1e12, 1)}
                node [pos = 0.12, anchor = south] {$M_\odot$};
            \addplot [only marks, mark = *, mark size = 1.5pt] coordinates {(1e9, 23.1)};
        \end{loglogaxis}
    \end{tikzpicture}
    \caption{Масса Джинса для молекулярного водорода при разных температурах; точка~--- облако из примера}
    \label{pic:jeans-mass}
\end{wrapfigure}
\paragraph{Пример} Звёзды образуются в молекулярных облаках: температура $T = 10~\text{К}$, концентрация молекул $n = 10^9~\text{м}^{-3}$, газ состоит в основном из молекулярного водорода, $m = 2m_p$, откуда $\rho = nm = 3.3 \cdot 10^{-18}~\text{кг}/\text{м}^3$. Подставляя в \eqref{eq:jeans-mass} и \eqref{eq:jeans-radius}, получаем, \lookPicRef{pic:jeans-mass},
\begin{gather*}
    M_J \simeq 4.6 \cdot 10^{31}~\text{кг} \simeq 23 M_\odot,\\
    R_J \simeq 1.5 \cdot 10^{16}~\text{м} \simeq 0.5~\text{пк},
\end{gather*}
а время свободного падения составляет около $1.5$ миллиона лет. Для сравнения, в разреженном атомарном газе с $T = 100~\text{К}$, $n = 10^7~\text{м}^{-3}$ и $m = m_p$ масса Джинса равна $3 \cdot 10^4 M_\odot$, а радиус Джинса~--- $30~\text{пк}$. Именно поэтому звёзды рождаются не в тёплом межзвёздном газе, а в его самых холодных и плотных сгустках.

\paragraph{Фрагментация} Что происходит с облаком после того, как сжатие началось? При сжатии гравитация совершает работу, и выделяющееся тепло нагревает газ. Судьба этого тепла зависит от того, насколько облако прозрачно для собственного излучения, \lookSecRef{sec:optical-thickness}. Оптическая толщина облака по порядку величины равна $\tau \sim n \sigma R$, где $\sigma$~--- сечение взаимодействия фотона с частицей газа. Пока $\tau \ll 1$, фотон, испущенный в любой точке облака, свободно выходит наружу: облако прозрачно, тепло уносится излучением, и температура газа почти не меняется, а плотность растёт. Масса Джинса при этом падает как $\rho^{-1/2}$~--- облако сползает вправо и вниз по своей кривой на \picRef{pic:jeans-mass},~--- и части облака, которые раньше удерживались давлением, сами начинают сжиматься. Облако распадается на всё более мелкие сгустки~--- \term{фрагментирует}. Так из одного облака массой в тысячи солнечных рождается не одна гигантская звезда, а целое звёздное скопление, \lookSecRef{sec:open-clusters}.

При сжатии сгустка постоянной массы концентрация растёт как $n \propto R^{-3}$, поэтому его оптическая толщина увеличивается, $\tau \propto R^{-2}$. Рано или поздно $\tau$ становится порядка единицы: фотон поглощается раньше, чем успевает покинуть сгусток, и тепло остаётся внутри. С этого момента сжатие идёт с ростом температуры, а масса Джинса, пропорциональная $T^{3/2}$, снова увеличивается. Фрагментация прекращается, и такой сгусток, продолжая сжиматься уже как целое, становится \term{протозвездой}.
